% 原书页29下部至44章2前；逐页图像见 evidence/p29.png--p44.png.

\section{Hecke算子}\label{sec:modular-1-3}

\subsection{一般形式体系}\label{sec:modular-1-3-1}

群 $G$ 的两个子群 $\Gamma_1,\Gamma_2$ 称为\newterm{modular-commensurable}{可公度的}{commensurable}，如果 $\Gamma_1\cap\Gamma_2$ 在二者中均为有限指数子群.这在 $G$ 的子群之间给出等价关系（习题：证明传递性）.若对每个 $\alpha\in G$，$\Gamma$ 与 $\alpha\Gamma\alpha^{-1}$ 可公度，就称 $\Gamma$ 为 $G$ 的\newterm{modular-hecke-subgroup}{Hecke子群}{Hecke subgroup}.此时，与 $\Gamma$ 可公度的每个子群也都是Hecke子群.

\begin{example*}

在 $G=\mathrm{GL}_2(\mathbb Q)$ 中，$\mathrm{SL}_2(\mathbb Z)$ 的有限指数子群都是Hecke子群.只须证明 $\alpha\Gamma(1)\alpha^{-1}\cap\Gamma(1)$ 在 $\Gamma(1)$ 中指数有限.适当乘以标量后可设 $\alpha\in\mathrm M_2(\mathbb Z)$.令 $N=\det\alpha$，则 $N\alpha^{-1}\in\mathrm M_2(\mathbb Z)$（这是 $\alpha$ 的余子式矩阵的转置），故 $\alpha^{-1}(N\mathrm M_2(\mathbb Z))\alpha\subset\mathrm M_2(\mathbb Z)$，从而 $\alpha^{-1}\Gamma(N)\alpha\subset\mathrm M_2(\mathbb Z)$.后者为群，因而包含于 $\mathrm{GL}_2(\mathbb Z)$；又因 $\det(\alpha^{-1}\gamma\alpha)=\det\gamma$，它包含于 $\mathrm{SL}_2(\mathbb Z)$.最终得到 $\Gamma(N)\subset\alpha\Gamma(1)\alpha^{-1}\cap\Gamma(1)$.

\end{example*}

设 $\Gamma_1,\Gamma_2$ 为 $G$ 中两个可公度的Hecke子群.$\Gamma_2$ 在 $\Gamma_1\backslash G$ 上右作用的轨道均为有限集，因为有双射
\begin{align*}
(\alpha^{-1}\Gamma_1\alpha\cap\Gamma_2)\backslash\Gamma_2&\xrightarrow{\sim}\Gamma_1\backslash(\Gamma_1\alpha\Gamma_2),\\
(\alpha^{-1}\Gamma_1\alpha\cap\Gamma_2)\gamma&\longmapsto\Gamma_1\alpha\gamma.
\end{align*}
因此可将 $\mathbb Z[\Gamma_1\backslash G/\Gamma_2]$ 与不变量 $\mathbb Z[\Gamma_1\backslash G]^{\Gamma_2}$ 等同.\footnote{对集合 $X$，$\mathbb Z[X]$ 表示以 $\{[x],x\in X\}$ 为基的自由 $\mathbb Z$ 模.}考虑映射
\begin{align*}
\operatorname{Hom}_{\mathbb Z[G]}(\mathbb Z[\Gamma_2\backslash G],\mathbb Z[\Gamma_1\backslash G])&\longrightarrow\mathbb Z[\Gamma_1\backslash G/\Gamma_2],&\varphi&\longmapsto\varphi([\Gamma_2.1]).
\end{align*}
这里Hom表示右 $\mathbb Z[G]$ 模同态.因为 $\mathbb Z[\Gamma_2\backslash G]$ 由 $[\Gamma_2.1]$ 生成为 $\mathbb Z[G]$ 模，同态完全由该元素的像确定.该像须为 $\Gamma_2$ 不变元，故可视为右端中的元素.反之，每个这样的元素都确定一个 $\varphi$，所以上述映射是双射.具体地，
\begin{align*}
[\Gamma_1\alpha\Gamma_2]:\mathbb Z[\Gamma_2\backslash G]&\longrightarrow\mathbb Z[\Gamma_1\backslash G],&[\Gamma_2g]&\longmapsto\sum_{\Gamma_1\gamma\in\Gamma_1\backslash(\Gamma_1\alpha\Gamma_2)}[\Gamma_1\gamma g].
\end{align*}
求和下标表示可令 $\gamma$ 遍历商集 $\Gamma_1\backslash(\Gamma_1\alpha\Gamma_2)$ 在 $G$ 中的任意代表系 $\mathcal L(\Gamma_1\alpha\Gamma_2)$；结果显然与选择无关.

若 $\Gamma_3$ 是第三个与 $\Gamma_1$（从而与 $\Gamma_2$）可公度的子群，同态复合给出乘积映射
\begin{align*}
\mathbb Z[\Gamma_1\backslash G/\Gamma_2]\otimes\mathbb Z[\Gamma_2\backslash G/\Gamma_3]&\longrightarrow\mathbb Z[\Gamma_1\backslash G/\Gamma_3],\\
[\Gamma_1\alpha\Gamma_2][\Gamma_2\beta\Gamma_3]&=\sum_\delta c_\delta[\Gamma_1\delta\Gamma_3],\\
c_\delta&=\big|\{(\Gamma_1\gamma,\gamma')\in\Gamma_1\backslash(\Gamma_1\alpha\Gamma_2)\times\mathcal L(\Gamma_2\beta\Gamma_3):\Gamma_1\gamma\gamma'=\Gamma_1\delta\}\big|.
\end{align*}
其中 $\mathcal L(\Gamma_2\beta\Gamma_3)$ 是相应左商集在 $G$ 中的代表系.这个集合依赖代表系的选择，但其基数不依赖，因为它正是上述乘积中 $[\Gamma_1\delta\Gamma_3]$ 的系数.［习题：直接验证这一基数与代表系无关.］

\begin{remark*}

$c_\delta\ne0$ 当且仅当 $\Gamma_1\delta\Gamma_3\subset\Gamma_1\alpha\Gamma_2\beta\Gamma_3$.

\end{remark*}

取 $\Gamma_1=\Gamma_2=\Gamma_3=:\Gamma$，便在 $\mathbb Z[\Gamma\backslash G/\Gamma]$ 上得到环结构，称为 $(G,\Gamma)$ 的\newterm{modular-hecke-ring}{Hecke环}{Hecke ring}，并得到环同构
\begin{align*}
\operatorname{End}_{\mathbb Z[G]}(\mathbb Z[\Gamma\backslash G])\xrightarrow{\sim}\mathbb Z[\Gamma\backslash G/\Gamma].
\end{align*}
若 $\Delta\subset G$ 对乘法封闭且包含 $\Gamma$，记 $\mathbb Z[\Gamma\backslash\Delta/\Gamma]$ 为所有 $[\Gamma\delta\Gamma]$（$\delta\in\Delta$）生成的子模.由乘积公式可知它是子环.

现在给定带有 $G$ 的线性右作用 $(v,g)\mapsto v.g$ 的复向量空间（或 $\mathbb Z$ 模）$V$.同理，映射
\begin{align*}
\operatorname{Hom}_{\mathbb Z[G]}(\mathbb Z[\Gamma\backslash G],V)&\longrightarrow V^\Gamma,&\varphi&\longmapsto\varphi([\Gamma.1])
\end{align*}
是复向量空间同构.与 $[\Gamma_1\alpha\Gamma_2]$ 复合给出线性映射 $V^{\Gamma_1}\to V^{\Gamma_2}$，具体为
\begin{align*}
\forall v\in V^{\Gamma_1},\qquad v.[\Gamma_1\alpha\Gamma_2]=\sum_{\Gamma_1\gamma\in\Gamma_1\backslash(\Gamma_1\alpha\Gamma_2)}v\gamma.
\end{align*}
［习题：直接验证右端在 $\Gamma_2$ 下不变.］按定义，这些映射与乘积相容；特别地，$\Gamma_1=\Gamma_2=\Gamma$ 时，$V^\Gamma$ 成为Hecke环上的右模.

\begin{example*}

取 $V$ 为函数 $\mathbb H\to\mathbb C$ 的空间，$G=\mathrm{GL}_2(\mathbb Q)_+$ 按权 $k\ge0$ 作用：
\begin{align*}
f[\gamma]_k(z)=\det(\gamma)^{k/2}j_\gamma(z)^{-k}f(\gamma z).
\end{align*}
$V^\Gamma$ 包含子空间 $\mathcal M_k(\Gamma),\mathcal S_k(\Gamma)$.公式
\begin{align*}
f.[\Gamma_1\alpha\Gamma_2]_k=\sum_{\gamma\in\mathcal L(\Gamma_1\alpha\Gamma_2)}f[\gamma]_k
\end{align*}
表明（习题）该作用保持 $\mathbb H$ 上及尖点处的全纯性，因而将 $\mathcal M_k(\Gamma_1)$、$\mathcal S_k(\Gamma_1)$ 分别映入 $\mathcal M_k(\Gamma_2)$、$\mathcal S_k(\Gamma_2)$.特别地，Hecke环右作用于 $\mathcal M_k(\Gamma)$ 和 $\mathcal S_k(\Gamma)$.

\end{example*}

\subsection{级为 \texorpdfstring{$\Gamma(1)$}{Γ(1)} 的Hecke环}\label{sec:modular-1-3-2}

令 $\Delta=\Delta(1)=\mathrm M_2(\mathbb Z)\cap\mathrm{GL}_2(\mathbb Q)_+$.显然它对乘法封闭且包含 $\Gamma(1)$.下面描述环
\begin{align*}
\mathbb Z[\Gamma(1)\backslash\Delta/\Gamma(1)].
\end{align*}
由初等因子定理，$\alpha\in\Delta$ 的余核 $\operatorname{Coker}(\alpha)=\mathbb Z^2/\alpha(\mathbb Z^2)$ 形如 $\mathbb Z/l\mathbb Z\oplus\mathbb Z/m\mathbb Z$，若要求 $l\mid m$，则 $l,m\ge1$ 唯一确定.

\begin{lemma*}

记 $\Delta_{l,m}$ 为初等因子为 $(l,m)$ 的矩阵集合.则

\begin{enumerate}[label=\roman*)]
\item
\begin{align*}
\Delta_{l,m}&=\Gamma(1)\begin{bmatrix}l&0\\0&m\end{bmatrix}\Gamma(1),
\end{align*}
\item
\begin{align*}
\mathcal L\left(\Gamma(1)\begin{bmatrix}l&0\\0&m\end{bmatrix}\Gamma(1)\right)&:=\left\{\begin{bmatrix}a&b\\0&d\end{bmatrix}:ad=lm,\ 0\le b<d,\ (a,b,d)=l\right\}
\end{align*}
是 $\Delta_{l,m}$ 中 $\Gamma(1)$ 左陪集的代表系.
\end{enumerate}

\end{lemma*}

\begin{proof}

i) 若 $\alpha=\gamma_1\left[\begin{smallmatrix}l&0\\0&m\end{smallmatrix}\right]\gamma_2$，$\gamma_i\in\Gamma(1)$，则 $\alpha(\mathbb Z^2)=\gamma_1(l\mathbb Z\oplus m\mathbb Z)$，所以 $\gamma_1$ 诱导同构 $\mathbb Z^2/(l\mathbb Z\oplus m\mathbb Z)\simeq\operatorname{Coker}\alpha$.

反之，若 $\alpha$ 的初等因子为 $(l,m)$，存在 $\mathbb Z^2$ 的基 $(\omega_1,\omega_2)$，使 $(l\omega_1,m\omega_2)$ 为 $\alpha(\mathbb Z^2)$ 的基.令 $\gamma_1$ 为标准基到 $(\omega_1,\omega_2)$ 的过渡矩阵，$\gamma_2$ 为 $(l\omega_1,m\omega_2)$ 到 $(\alpha(e_1),\alpha(e_2))$ 的过渡矩阵，便有 $\alpha=\gamma_1\left[\begin{smallmatrix}l&0\\0&m\end{smallmatrix}\right]\gamma_2$.起初只知 $\gamma_i\in\mathrm{GL}_2(\mathbb Z)$；必要时分别用 $\gamma_1\left[\begin{smallmatrix}1&0\\0&-1\end{smallmatrix}\right]$ 和 $\left[\begin{smallmatrix}1&0\\0&-1\end{smallmatrix}\right]\gamma_2$ 替换，即得结论.

ii) 使用以下观察：对 $\beta,\beta'\in\Delta$，
\begin{align*}
\beta\Gamma(1)=\beta'\Gamma(1)\quad\Longleftrightarrow\quad\beta(\mathbb Z^2)=\beta'(\mathbb Z^2).
\end{align*}
正向显然；反向取两组像基之间的过渡矩阵 $\gamma$，即得 $\beta'=\beta\gamma$.但我们需要左陪集，故对转置矩阵作此论证.

设 $\alpha\in\Delta_{l,m}$，$\Lambda={}^t\alpha(\mathbb Z^2)$.上述观察说明，$\left[\begin{smallmatrix}a&b\\0&d\end{smallmatrix}\right]\in\Gamma(1)\alpha$ 当且仅当 $\Lambda$ 由 $(ae_1+be_2,de_2)$ 生成.此时 $d$ 是使 $de_2\in\Lambda$ 的最小正整数，$a$ 由 $ad=lm$ 确定，$b$ 模 $d$ 也随之确定，故满足 $0\le b<d$ 的矩阵至多一个.为证明存在性，考虑正合列
\[
\begin{tikzcd}[column sep=small]
0 \arrow[r] & (\Lambda+\mathbb Ze_2)/\Lambda \arrow[r] & \mathbb Z^2/\Lambda \arrow[r] & \mathbb Z^2/(\Lambda+\mathbb Ze_2) \arrow[r] & 0
\end{tikzcd}
\]
左项是由 $e_2$ 的像生成的循环群，记其阶为 $d$；右项由 $e_1$ 的像生成，记其阶为 $a$.于是 $ad=|\mathbb Z^2/\Lambda|=|\det({}^t\alpha)|=lm$，且 $(ae_1+\mathbb Ze_2)\cap\Lambda\ne\varnothing$.对其中任一 $\omega_1$，$(\omega_1,de_2)$ 是 $\Lambda$ 的基；因 $de_2\in\Lambda$，可取 $\omega_1=ae_1+be_2$，$0\le b<d$.还须证明 $(a,b,d)=l$.转置仍属 $\Delta_{l,m}$，故 $\mathbb Z^2/\Lambda\simeq\mathbb Z/l\mathbb Z\oplus\mathbb Z/m\mathbb Z$.其乘 $l$ 映射的核同构于 $(\mathbb Z/l\mathbb Z)^2$，这等价于 $l\mathbb Z^2\supset\Lambda$，从而 $l\mid(a,b,d)$.反之，令 $\lambda=(a,b,d)$，则 $\Lambda\subset\lambda\mathbb Z^2$，$\mathbb Z^2/\Lambda$ 包含 $(\mathbb Z/\lambda\mathbb Z)^2$，因而 $\lambda\mid l$.

\end{proof}

记 $T(l,m)=[\Gamma(1)\left[\begin{smallmatrix}l&0\\0&m\end{smallmatrix}\right]\Gamma(1)]$.这些元素（$1\le l\mid m$）组成Hecke环的一组基，乘法表如下.

\begin{proposition*}

对素数 $p$ 有

\begin{enumerate}[label=\roman*)]
\item
\begin{align*}
T(l,m)&=T(l,l)T(1,m/l)=T(1,m/l)T(l,l),
\end{align*}
\item
\begin{align*}
T(l,m)T(l',m')&=T(ll',mm')\quad((lm,l'm')=1),
\end{align*}
\item
\begin{align*}
T(1,p)T(1,p^r)&=T(1,p^{r+1})+\begin{cases}(p+1)T(p,p),&r=1,\\pT(p,p)T(1,p^{r-1}),&r>1.\end{cases}
\end{align*}
\end{enumerate}

特别地，该Hecke环是以所有 $T(1,p),T(p,p)$ 为变元的多项式环.

\end{proposition*}

\begin{proof}

i) 更一般地，若 $z$ 在 $G$ 的中心，则乘法公式给出 $[\Gamma z\Gamma][\Gamma\alpha\Gamma]=[\Gamma z\alpha\Gamma]=[\Gamma\alpha\Gamma][\Gamma z\Gamma]$.

ii) 由 i) 可设 $l=l'=1$.先证明集合等式
\begin{align*}
\Gamma(1)\begin{bmatrix}1&0\\0&m\end{bmatrix}\Gamma(1)\begin{bmatrix}1&0\\0&m'\end{bmatrix}\Gamma(1)=\Gamma(1)\begin{bmatrix}1&0\\0&mm'\end{bmatrix}\Gamma(1).
\end{align*}
由引理，只须证明初等因子分别为 $(1,m),(1,m')$ 的 $\alpha,\alpha'$ 的乘积具有初等因子 $(1,mm')$.正合列
\[
\begin{tikzcd}[column sep=small]
0 \arrow[r] & \alpha\mathbb Z^2/\alpha\alpha'\mathbb Z^2 \arrow[r] & \mathbb Z^2/\alpha\alpha'\mathbb Z^2 \arrow[r] & \mathbb Z^2/\alpha\mathbb Z^2 \arrow[r] & 0
\end{tikzcd}
\]
表明余核是 $\mathbb Z/m\mathbb Z$ 被 $\mathbb Z/m'\mathbb Z$ 的Abel扩张.由 $(m,m')=1$，该扩张同构于 $\mathbb Z/m\mathbb Z\oplus\mathbb Z/m'\mathbb Z=\mathbb Z/mm'\mathbb Z$，故结论成立.于是 $T(1,m)T(1,m')=cT(1,mm')$，其中
\begin{align*}
C=\{(\alpha,\alpha')\in\mathcal L(\Delta_{1,m})\times\mathcal L(\Delta_{1,m'}):\Gamma(1)\alpha\alpha'=\Gamma(1)\begin{bmatrix}1&0\\0&mm'\end{bmatrix}\},\quad c=|C|.
\end{align*}
使用引理的代表系，写 $\alpha=\left[\begin{smallmatrix}a&b\\0&d\end{smallmatrix}\right]$，$ad=m$，$0\le b<d$，$(a,b,d)=1$；类似地写 $\alpha'=\left[\begin{smallmatrix}a'&b'\\0&d'\end{smallmatrix}\right]$，$a'b'=m'$，$0\le b'<d'$，$(a',b',d')=1$.\footnote{译注：原书页33此处确写 $a'b'=m'$，与对角元乘积的条件不一致；保留原式.}则
\begin{align*}
\alpha\alpha'=\begin{bmatrix}aa'&ab'+bd'\\0&dd'\end{bmatrix}.
\end{align*}
由引理证明，这属于指定左陪集当且仅当 $aa'=1$，$dd'=mm'$，$mm'\mid ab'+bd'$.这又等价于 $a=a'=1$，$d=m$，$d'=m'$（因 $m,m'$ 互素），且 $mm'\mid b'+bm'\Longleftrightarrow b'+bm'=0\Longleftrightarrow b=b'=0$.所以 $C$ 只有对角矩阵组成的那一对，$c=1$.

iii) 同一正合列表明，初等因子为 $(1,p),(1,p^r)$ 的矩阵乘积的余核，是 $\mathbb Z/p\mathbb Z$ 被 $\mathbb Z/p^r\mathbb Z$ 的扩张.因此它或为 $\mathbb Z/p^{r+1}\mathbb Z$，或为 $\mathbb Z/p\mathbb Z\oplus\mathbb Z/p^r\mathbb Z$，初等因子分别为 $(1,p^{r+1})$、$(p,p^r)$.所以相应双陪集的集合乘积是 $\Delta_{1,p^{r+1}}\sqcup\Delta_{p,p^r}$，而
\begin{align*}
T(1,p)T(1,p^r)=cT(1,p^{r+1})+c'T(p,p^r).
\end{align*}
$c$ 是集合 $C=\{(\alpha,\alpha')\in\mathcal L(\Delta_{1,p})\times\mathcal L(\Delta_{1,p^r}):\Gamma(1)\alpha\alpha'=\Gamma(1)\left[\begin{smallmatrix}1&0\\0&p^{r+1}\end{smallmatrix}\right]\}$ 的基数.引理给出的代表系为
\begin{align*}
\mathcal L(\Delta_{1,p^r})=\left\{\begin{bmatrix}p^{r-s}&m\\0&p^s\end{bmatrix}:0\le s\le r,\ m<p^s,\ (m,p^s,p^{r-s})=1\right\}.
\end{align*}
于是 $C=\{(\left[\begin{smallmatrix}1&0\\0&p\end{smallmatrix}\right],\left[\begin{smallmatrix}1&0\\0&p^r\end{smallmatrix}\right])\}$，故 $c=1$.\footnote{译注：原书页34这份代表系的条件只写 $m<p^s$，未重写此前已有的非负条件；此处保留.}

也可以直接列出 $C'=\{(\alpha,\alpha')\in\mathcal L(\Delta_{1,p})\times\mathcal L(\Delta_{1,p^r}):\Gamma(1)\alpha\alpha'=\Gamma(1)\left[\begin{smallmatrix}p&0\\0&p^r\end{smallmatrix}\right]\}$ 来求 $c'$；但使用双陪集的次数更简洁.定义 $\deg([\Gamma\alpha\Gamma])=|\Gamma\backslash(\Gamma\alpha\Gamma)|$，再线性延拓到Hecke环.可验证（见下面习题）这是环同态.在这里 $\deg\Delta_{1,p^r}=p^r+p^{r-1}$，故
\begin{align*}
(1+p)(p^r+p^{r-1})&=c(p^{r+1}+p^r)+c'\deg\Delta_{p,p^r},\\
\deg\Delta_{p,p^r}=\deg\Delta_{1,p^{r-1}}&=\begin{cases}1,&r=1,\\p^{r-1}+p^{r-2},&r>1.\end{cases}
\end{align*}
由 $c=1$ 得 $c'=1+p$（$r=1$）或 $c'=p$（$r>1$）.

最后，由 i)、ii)，所有 $T(1,p),T(p',p')$ 彼此交换；由 iii)、ii)，它们生成Hecke环.还须验证代数独立性.令
\begin{align*}
P(l,m)=\prod_{p\mid l}T(p,p)^{v_p(l)}\prod_{p\mid m/l}T(1,p)^{v_p(m/l)}.
\end{align*}
对 $(l,m)$ 取任意细化由 $m\mid m'$ 诱导的偏序的严格次序，ii)、iii) 表明把 $P(l,m)$ 用 $T(l,m)$ 表出的矩阵是对角元为1的上三角矩阵.后者为一组基，故前者也为一组基.

\end{proof}

\begin{exercise*}

证明在 $[\Gamma\alpha\Gamma][\Gamma\beta\Gamma]=\sum_\delta c_\delta[\Gamma\delta\Gamma]$ 中，
\begin{align*}
c_\delta=\frac{|\{(\gamma,\gamma')\in\mathcal L(\Gamma\alpha\Gamma)\times\mathcal L(\Gamma\beta\Gamma):\Gamma\gamma\gamma'\Gamma=\Gamma\delta\Gamma\}|}{|\Gamma\backslash(\Gamma\delta\Gamma)|}.
\end{align*}
由此证明次数映射确为环同态.还要证明，$[\Gamma\alpha\Gamma][\Gamma\beta\Gamma]=[\Gamma\alpha\beta\Gamma]$ 等价于 $\mathcal L(\Gamma\alpha\Gamma)\mathcal L(\Gamma\beta\Gamma)$（$\Delta$ 中的乘积）是 $\Gamma\backslash(\Gamma\alpha\beta\Gamma)$ 的代表系.

\end{exercise*}

\begin{corollary*}

对 $n\in\mathbb N$，记 $T(n)=\sum_{lm=n}T(l,m)$，特别地 $T(p)=T(1,p)$.则

\begin{enumerate}[label=\roman*)]
\item
\begin{align*}
T(nm)&=T(n)T(m)\quad((n,m)=1),
\end{align*}
\item
\begin{align*}
T(p^{r+1})&=T(p)T(p^r)-pT(p,p)T(p^{r-1}).
\end{align*}
\end{enumerate}

等价地，$T(n)T(m)=\sum_{l\mid(m,n)}lT(l,l)T(mn/l^2)$.

\end{corollary*}

\begin{proof}留作习题.\end{proof}

\subsection{在级为1的模形式上的作用}\label{sec:modular-1-3-3}

前面已说明，Hecke环在 $\Gamma(1)$ 的权 $k$ 自守函数上的作用保持 $\mathcal M_k(\Gamma(1))$ 和 $\mathcal S_k(\Gamma(1))$.此作用由 $G=\mathrm{GL}_2(\mathbb Q)_+$ 的权 $k$ 作用 $f[\gamma]_k(z)=\det(\gamma)^{k/2}j_\gamma(z)^{-k}f(\gamma z)$ 诱导.不过，为使公式具有有理乃至整数系数，最好改用归一化
\begin{align*}
f[\gamma]'_k(z)=\det(\gamma)^{k-1}j_\gamma(z)^{-k}f(\gamma z)=\det(\gamma)^{k/2-1}f[\gamma]_k(z).
\end{align*}
相应地，$[\Gamma(1)\alpha\Gamma(1)]'_k=\det(\alpha)^{k/2-1}[\Gamma(1)\alpha\Gamma(1)]_k$.具体公式如下.

\begin{lemma*}

对 $n\ge1$ 和 $f\in\mathcal M_k(\Gamma(1))$，
\begin{align*}
f[T(n)]'_k(z)=n^{k-1}\sum_{ad=n}\sum_{b=0}^{d-1}d^{-k}f((az+b)/d).
\end{align*}
若 $f(z)=\sum_{m\in\mathbb N}a(m)q^m$ 是其在 $\infty$ 的 $q$ 展开（$q=\exp(2i\pi z)$），则
\begin{align*}
f[T(n)]'_k(z)=\sum_{m\in\mathbb N}\left(\sum_{d\mid(m,n)}d^{k-1}a(mn/d^2)\right)q^m.
\end{align*}
\end{lemma*}

\begin{proof}

按定义，对 $\Delta_{l,m}$ 中任意左陪集代表系 $\mathcal L$，有
\begin{align*}
f[T(l,m)]'_k(z)=\sum_{\gamma\in\mathcal L}\det(\gamma)^{k-1}j_\gamma(z)^{-k}f(\gamma z).
\end{align*}
取上个引理的代表系
\begin{align*}
\mathcal L=\left\{\begin{bmatrix}a&b\\0&d\end{bmatrix}:ad=lm,\ 0\le b<d,\ (a,b,d)=l\right\},
\end{align*}
便得与所述公式类似、但限于 $(a,b,d)=l$ 的和.再对 $lm=n$ 的 $(l,m)$ 求和，即得第一式.第二式由第一式以及
\begin{align*}
f((az+b)/d)&=\sum_{m\in\mathbb N}a(m)e^{2i\pi mb/d}q^{ma/d},&\sum_{b=0}^{d-1}e^{2i\pi mb/d}&=\begin{cases}d,&d\mid m,\\0,&\text{否则}\end{cases}
\end{align*}
得出.

\end{proof}

下面证明 $T(n)$ 关于Petersson内积自伴.先给出一个引理.

\begin{lemma*}

设 $\alpha\in\Delta$，$\Gamma$ 是 $\Gamma(1)\cap\alpha^{-1}\Gamma(1)\alpha$ 的有限指数子群.对 $f,g\in\mathcal S_k(\alpha\Gamma\alpha^{-1})$，有 $f[\alpha]_k,g[\alpha]_k\in\mathcal S_k(\Gamma)$，且
\begin{align*}
\langle f[\alpha]_k,g[\alpha]_k\rangle_\Gamma=\langle f,g\rangle_{\alpha\Gamma\alpha^{-1}}.
\end{align*}
\end{lemma*}

\begin{proof}

前一个断言显然.令 $h(z)=\Im(z)^kf(z)\overline{g(z)}$.由 $\Im(\alpha z)=\det(\alpha)|j_\alpha(z)|^{-2}\Im(z)$ 得 $h(\alpha z)=\Im(z)^kf[\alpha]_k(z)\overline{g[\alpha]_k(z)}$，于是
\begin{align*}
\langle f[\alpha]_k,g[\alpha]_k\rangle_\Gamma
&=\frac1{V_\Gamma}\int_{D_\Gamma}h(\alpha z)d\mu(z)
=\frac1{V_\Gamma}\int_{\alpha D_\Gamma}h(z)d\mu(z)
=\langle f,g\rangle_{\alpha\Gamma\alpha^{-1}}.
\end{align*}
这里 $D_\Gamma$ 是基本域.第一步用内积定义，第二步用 $d\mu$ 在 $\mathrm{GL}_2(\mathbb R)_+$ 下的不变性，第三步用 $\alpha D_\Gamma$ 是共轭群的基本域，且体积与 $D_\Gamma$ 相同.

\end{proof}

\begin{proposition*}

对 $f,g\in\mathcal S_k(\Gamma(1))$，$\langle f[T(n)]'_k,g\rangle_{\Gamma(1)}=\langle f,g[T(n)]'_k\rangle_{\Gamma(1)}$.

\end{proposition*}

\begin{proof}

只须对 $[T(l,m)]_k$ 证明同一公式.先找出 $\Delta_{l,m}$ 中左右陪集共用的代表系.取左陪集代表系 $\mathcal L$；转置保持 $\Delta_{l,m}$，故 ${}^t\mathcal L$ 是右陪集代表系.对 $\alpha\in\mathcal L$，取 $\gamma_\alpha,\delta_\alpha\in\Gamma(1)$ 使 ${}^t\alpha=\gamma_\alpha\alpha\delta_\alpha$.令 $\widetilde\alpha=\gamma_\alpha\alpha={}^t\alpha\delta_\alpha^{-1}$，则 $\Gamma(1)\widetilde\alpha=\Gamma(1)\alpha$，$\widetilde\alpha\Gamma(1)={}^t\alpha\Gamma(1)$.故 $\widetilde{\mathcal L}=\{\widetilde\alpha:\alpha\in\mathcal L\}$ 满足要求.

考虑反对合 $\alpha^*=\det(\alpha)\alpha^{-1}$.它保持 $\Gamma(1)$，并把 $\left[\begin{smallmatrix}l&0\\0&m\end{smallmatrix}\right]$ 映到 $\left[\begin{smallmatrix}m&0\\0&l\end{smallmatrix}\right]$，故保持 $\Delta_{l,m}$.因此 $\widetilde{\mathcal L}^*$ 是左陪集代表系.又对任意 $\alpha\in\Delta$，$f[\alpha^{-1}]_k=f[\alpha^*]_k$.

选取 $\Gamma(1)$ 的正规子群 $\Gamma$，使它包含于 $\Gamma(1)\cap\left[\begin{smallmatrix}l&0\\0&m\end{smallmatrix}\right]^{-1}\Gamma(1)\left[\begin{smallmatrix}l&0\\0&m\end{smallmatrix}\right]$.于是对所有 $\alpha\in\widetilde{\mathcal L}$，$\Gamma$ 都包含于 $\alpha\Gamma(1)\alpha^{-1}$ 和 $\alpha^{-1}\Gamma(1)\alpha$.由引理，
\begin{align*}
\langle f[T(l,m)]_k,g\rangle_{\Gamma(1)}&=\langle f[T(l,m)]_k,g\rangle_\Gamma\\
&=\sum_{\alpha\in\widetilde{\mathcal L}}\langle f[\alpha]_k,g\rangle_\Gamma
=\sum_{\alpha\in\widetilde{\mathcal L}}\langle f[\alpha]_k,g\rangle_{\Gamma(1)\cap\alpha^{-1}\Gamma(1)\alpha}\\
&=\sum_{\alpha\in\widetilde{\mathcal L}}\langle f,g[\alpha^{-1}]_k\rangle_{\alpha\Gamma(1)\alpha^{-1}\cap\Gamma(1)}
=\sum_{\alpha\in\widetilde{\mathcal L}}\langle f,g[\alpha^*]_k\rangle_\Gamma\\
&=\langle f,g[T(l,m)]_k\rangle_\Gamma=\langle f,g[T(l,m)]_k\rangle_{\Gamma(1)}.
\end{align*}
\end{proof}

\begin{corollary*}

$\mathcal S_k(\Gamma(1))$ 有一组正交基，其元素是所有 $T(n)$（乃至所有 $T(l,m)$）的共同特征向量.若 $f(z)=\sum_{n>0}a(n)q^n$ 是特征值为 $(\lambda(n))_{n\in\mathbb N^*}$ 的特征形式，则 $a(n)=a(1)\lambda(n)$（$n\in\mathbb N^*$）.特别地，若归一化为 $a(1)=1$，则
\begin{align*}
a(nm)&=a(n)a(m)\quad((n,m)=1),\\
a(p^{r+1})&=a(p)a(p^r)-p^{k-1}a(p^{r-1})\quad(p\text{ 为素数},\ r\ge1).
\end{align*}
\end{corollary*}

\begin{proof}

自伴算子 $T(l,m)$ 可对角化，而Hecke环交换，故能同时对角化，得到所述正交基.由上面的展开公式，$f[T(n)]'_k$ 的 $q$ 项系数为 $a(n)$，$\lambda(n)f$ 的相应系数为 $a(1)\lambda(n)$，故二者相等.当 $a(1)=1$ 时，所述关系等价于特征值的相同关系.$\lambda(n)$ 是代数同态 $\lambda:\mathbb Z[\Gamma(1)\backslash\Delta/\Gamma(1)]\to\mathbb C$ 在 $T(n)$ 处的值，故
\begin{align*}
\lambda(nm)&=\lambda(n)\lambda(m)\quad((n,m)=1),\\
\lambda(p^{r+1})&=\lambda(p)\lambda(p^r)-p\lambda(T(p,p))a(p^{r-1})\quad(r\ge1).
\end{align*}
只须计算 $\lambda(T(p,p))$.原式为
\begin{align*}
f[T(p,p)]'_k(z)=\det\begin{bmatrix}p&0\\0&p\end{bmatrix}^{k-1}j_{\left[\begin{smallmatrix}p&0\\0&p\end{smallmatrix}\right]}(z)^{-1}f(z)=p^{2k-2}p^{-k}f(z),
\end{align*}
所以 $\lambda(T(p,p))=p^{k-2}$.\footnote{译注：原书页37的 $j$ 指数确为 $-1$，而后面的计算使用 $-k$；保留原式.}

\end{proof}

\begin{remark*}

$a(n)=\lambda(n)a(1)$ 表明，尖点特征形式（与任何特征向量一样非零）必有 $a(1)\ne0$，所以它是某个归一化特征形式的倍数.

\end{remark*}

\begin{example*}

$\mathcal S_{12}(\Gamma(1))$ 的维数为1，所以 $\Delta$ 是Hecke算子的特征向量.从而得到Ramanujan所猜想的 $\tau$ 函数的乘法性质.

\end{example*}

\begin{remark*}

Ramanujan关于 $\tau(n)$ 大小的第二个猜想可推广如下：若 $f\in\mathcal S_k(\Gamma(1))$ 是归一化特征形式，则对每个 $\varepsilon>0$，有 $a(n)=o(n^{(k-1)/2+\varepsilon})$.这一Ramanujan--Petersson猜想等价于（见习题课）$X^2-a(p)X+p^{k-1}$ 的根 $u_p,v_p$ 互为复共轭（因而模均为 $p^{(k-1)/2}$）.Shimura的工作说明，它可由关于有限域上某些簇有理点数大小的Weil猜想推出；后者最终由Deligne借助Grothendieck的平展上同调证明.

\end{remark*}

\begin{remark*}

（见习题课.）对归一化尖点特征形式，Dirichlet级数 $L(f,s)=\sum a(n)n^{-s}$ 在 $\Re(s)>k/2$ 收敛，并有Euler乘积
\begin{align*}
L(s,f)=\prod_p\frac1{1-a(p)p^{-s}+p^{k-1-2s}}.
\end{align*}
Ramanujan--Petersson猜想蕴含它在 $\Re(s)>(k+1)/2$ 收敛.\footnote{译注：原书页37先写 $\Re(s)>k/2$，末句又给出 $\Re(s)>(k+1)/2$；这里照录，不调整收敛范围.}

\end{remark*}

\subsection{级为 \texorpdfstring{$\Gamma_0(N)$}{Γ₀(N)} 的Hecke环}\label{sec:modular-1-3-4}

下面描述 $\mathcal H_0(N)=\mathbb Z[\Gamma_0(N)\backslash\Delta_0(N)/\Gamma_0(N)]$，其中
\begin{align*}
\Delta_0(N)=\left\{\begin{bmatrix}a&b\\c&d\end{bmatrix}\in\Delta:c\equiv0\pmod N,\ (a,N)=1\right\}.
\end{align*}
$\alpha\in\Delta_0(N)$ 的主要不变量仍是余核的初等因子 $(l,m)$；条件迫使 $(l,N)=1$.

\begin{lemma*}

记 $\Delta_0(N)_{l,m}$ 为不变量是 $(l,m)$ 的矩阵集，则

\begin{enumerate}[label=\roman*)]
\item
\begin{align*}
\Delta_0(N)_{l,m}&=\Gamma_0(N)\begin{bmatrix}l&0\\0&m\end{bmatrix}\Gamma_0(N),
\end{align*}
\item
\begin{align*}
\mathcal L_0(N)_{l,m}&=\left\{\begin{bmatrix}a&b\\0&d\end{bmatrix}:(a,N)=1,\ ad=lm,\ 0\le b<d,\ (a,b,d)=l\right\}
\end{align*}
是其中 $\Gamma_0(N)$ 左陪集的代表系.
\end{enumerate}

\end{lemma*}

\begin{proof}

i) 与 $\Gamma(1)$ 情形的差别在于：对 $\alpha\in\Delta_0(N)$，须找到 $\mathbb Z^2$ 的基 $(\omega_1,\omega_2)$，使 $(l\omega_1,m\omega_2)$ 是 $\alpha(\mathbb Z^2)$ 的基，$\omega_1\in\mathbb Z\oplus N\mathbb Z$（保证 $\gamma_1\in\Gamma_0(N)$），且 $\alpha(e_1)\in l\mathbb Z\omega_1\oplus Nm\mathbb Z\omega_2$（保证 $\gamma_2\in\Gamma_0(N)$）.可证，只要 $(l\omega_1,mN\omega_2)$ 是 $\alpha(\mathbb Z\oplus N\mathbb Z)$ 的基，就满足要求（习题，或见Miake 4.5.1）.

ii) 仿照 $\Gamma(1)$ 情形；注意若 $\alpha,\alpha'\in\Delta_0(N)$，且 $\alpha'=\gamma\alpha$，$\gamma\in\Gamma(1)$，则 $\gamma\in\Gamma_0(N)$.

\end{proof}

仍记 $T(l,m)=[\Gamma_0(N)\left[\begin{smallmatrix}l&0\\0&m\end{smallmatrix}\right]\Gamma_0(N)]$.由引理，这些元素（$1\le l\mid m$，$(l,N)=1$）组成 $\mathcal H_0(N)$ 的一组基.

\begin{proposition*}

对素数 $p$ 有

\begin{enumerate}[label=\roman*)]
\item
\begin{align*}
T(l,m)&=T(l,l)T(1,m/l)=T(1,m/l)T(l,l),
\end{align*}
\item
\begin{align*}
T(l,m)T(l',m')&=T(ll',mm')\quad((lm,l'm')=1),
\end{align*}
\item
\begin{align*}
T(1,p)T(1,p^r)&=T(1,p^{r+1})+\begin{cases}(p+1)T(p,p),&r=1,\\pT(p,p)T(1,p^{r-1}),&r>1\end{cases}\quad(p\nmid N),
\end{align*}
\item
\begin{align*}
T(1,p)T(1,p^r)&=T(1,p^{r+1})\quad(p\mid N).
\end{align*}
\end{enumerate}

特别地，Hecke环是以所有素数对应的 $T(1,p)$ 和所有不整除 $N$ 的素数对应的 $T(p',p')$ 为变元的多项式环.

\end{proposition*}

\begin{proof}

由引理，特别是代表系 $\mathcal L_0(N)_{l,m}$，可完全仿照 $\Gamma(1)$ 情形的论证.

\end{proof}

\begin{corollary*}

记 $T(n)=\sum_{lm=n,(l,N)=1}T(l,m)$，特别地 $T(p)=T(1,p)$.则

\begin{enumerate}[label=\roman*)]
\item
\begin{align*}
T(nm)&=T(n)T(m)\quad((n,m)=1),
\end{align*}
\item
\begin{align*}
T(p^{r+1})&=T(p)T(p^r)-pT(p,p)T(p^{r-1})\quad(p\nmid N),
\end{align*}
\item
\begin{align*}
T(p^{r+1})&=T(p)T(p^r)\quad(p\mid N).
\end{align*}
\end{enumerate}

等价地，$T(n)T(m)=\sum_{l\mid(m,n),(l,N)=1}lT(l,l)T(mn/l^2)$.

\end{corollary*}

\begin{proof}留作习题.\end{proof}

\subsection{在级为 \texorpdfstring{$\Gamma_1(N)$}{Γ₁(N)} 的模形式上的作用}\label{sec:modular-1-3-5}

当然，$\mathcal H_0(N)$ 自然作用于 $\mathcal M_k(\Gamma_0(N))$.实际上，把级为 $\Gamma_1(N)$ 的模形式适当分解后，它也作用于这些形式.出发点是 $\Gamma_1(N)$ 在 $\Gamma_0(N)$ 中正规，商群为有限Abel群.因此可分解为
\begin{align*}
\mathcal M_k(\Gamma_1(N))&=\bigoplus_\chi\mathcal M_k(N,\chi),&f&=\sum_\chi f_\chi,\\
f_\chi&=[\Gamma_0(N):\Gamma_1(N)]^{-1}\sum_{\gamma\in\Gamma_0(N)/\Gamma_1(N)}\chi^{-1}(\gamma)f[\gamma]_k.
\end{align*}
这里 $\chi$ 遍历特征标（即同态）$\Gamma_0(N)/\Gamma_1(N)\to\mathbb C^\times$，且
\begin{align*}
\mathcal M_k(N,\chi)=\{f\in\mathcal M_k(\Gamma_1(N)):f[\gamma]_k=\chi(\gamma)f,\ \forall\gamma\in\Gamma_0(N)\}.
\end{align*}
换言之，这是对由 $f[\gamma]_k^\chi=\chi^{-1}(\gamma)f[\gamma]_k$ 给出的权 $k$、以 $\chi$ 扭转的作用保持不变的函数空间.关键在于，$\chi$ 可延拓为乘法映射 $\Delta_0(N)\to\mathbb C^\times$，如下所述.因此可以相同公式 $f[\alpha]_k^\chi=\chi(\alpha)^{-1}f[\alpha]'_k$ 扭转 $\Delta_0(N)$ 的作用，并按通常公式定义Hecke环的作用：
\begin{align*}
f[\Gamma_0(N)\alpha\Gamma_0(N)]_k^\chi
&=\sum_{\alpha'\in\Gamma_0(N)\backslash(\Gamma_0(N)\alpha\Gamma_0(N))}f[\alpha']_k^\chi\\
&=\sum_{\alpha'\in\Gamma_0(N)\backslash(\Gamma_0(N)\alpha\Gamma_0(N))}\chi(\alpha')^{-1}f[\alpha']'_k.
\end{align*}
要说明延拓存在，注意 $\left[\begin{smallmatrix}a&b\\c&d\end{smallmatrix}\right]\mapsto d\pmod N$ 诱导同构 $\Gamma_0(N)/\Gamma_1(N)\simeq(\mathbb Z/N\mathbb Z)^\times$.故其特征标均形如
\begin{align*}
\chi\left(\begin{bmatrix}a&b\\c&d\end{bmatrix}\right)=\psi(d)=\psi(a)^{-1}=\overline\psi(a),
\end{align*}
其中 $\psi$ 遍历模 $N$ 的Dirichlet特征标.于是可用 $\overline\psi(a)$ 将它延拓到 $\Delta_0(N)$；却不能用 $\psi(d)$，因为 $d$ 未必与 $N$ 互素！

今后略为滥用记号，用同一字母 $\chi$ 表示 $\Gamma_0(N)$ 的特征标、对应的Dirichlet特征标，乃至对应的函数 $\mathbb Z\to\mathbb C$：在 $(n,N)=1$ 时取值 $\chi(n)$，其余情形取0.

\begin{lemma*}

对 $n\ge1$ 和 $f\in\mathcal M_k(N,\chi)$，
\begin{align*}
f[T(n)]_k^\chi(z)&=n^{k-1}\sum_{ad=n}\sum_{b=0}^{d-1}\chi(a)d^{-k}f((az+b)/d).
\end{align*}
若 $f(z)=\sum_{m\in\mathbb N}a(m)q^m$ 为它在 $\infty$ 的 $q$ 展开，$q=\exp(2i\pi z)$，则
\begin{align*}
f[T(n)]_k^\chi(z)=\sum_{m\in\mathbb N}\left(\sum_{d\mid(m,n)}\chi(d)d^{k-1}a(mn/d^2)\right)q^m.
\end{align*}
\end{lemma*}

\begin{proof}

与 $\Gamma(1)$ 情形作相同计算.在第一式中须记住 $(a,N)\ne1$ 时 $\chi(a)=0$；第二式的 $\chi(d)$ 亦然.

\end{proof}

接着考察 $T(n)$ 关于Petersson内积的伴随性质.这里与 $\Gamma(1)$ 情形有实质区别.先做一道习题.

\begin{exercise*}

验证分解 $\mathcal S_k(\Gamma_1(N))=\bigoplus_\chi\mathcal S_k(N,\chi)$ 关于Petersson内积是正交直和.

\end{exercise*}

\begin{proposition*}

设 $(n,N)=1$.对 $f,g\in\mathcal S_k(N,\chi)$，有
\begin{align*}
\langle f[T(n)]_k^\chi,g\rangle_{\Gamma_1(N)}=\overline\chi(n)\langle f,g[T(n)]_k^\chi\rangle_{\Gamma_1(N)}.
\end{align*}
特别地，$[T(n)]_k^\chi$ 是正规算子.

\end{proposition*}

\begin{proof}

只须在 $(lm,N)=1$ 时证明 $[T(l,m)]_k^\chi$ 的相应结论.仍用 $\alpha^*=\det(\alpha)\alpha^{-1}$.关键是此时 $\left[\begin{smallmatrix}m&0\\0&l\end{smallmatrix}\right]$ 仍在 $\Delta_0(N)$ 内，而 $\Gamma_0(N)$ 也被此对合保持，所以 $\Delta_0(N)_{l,m}$ 被保持.因此可找出左右陪集共用的代表系 $\widetilde{\mathcal L}$（此处不能用转置）.随后仿照前面计算，选取正规子群 $\Gamma$ 包含于 $\Gamma_1(N)\cap\left[\begin{smallmatrix}l&0\\0&m\end{smallmatrix}\right]^{-1}\Gamma_1(N)\left[\begin{smallmatrix}l&0\\0&m\end{smallmatrix}\right]$，则
\begin{align*}
\langle f[T(l,m)]_k^\chi,g\rangle_{\Gamma_1(N)}&=\langle f[T(l,m)]_k^\chi,g\rangle_\Gamma\\
&=\sum_{\alpha\in\widetilde{\mathcal L}}\langle f[\alpha]_k^\chi,g\rangle_\Gamma\\
&=\sum_{\alpha\in\widetilde{\mathcal L}}\chi^{-1}(\alpha)\langle f[\alpha]_k,g\rangle_{\Gamma_1(N)\cap\alpha^{-1}\Gamma_1(N)\alpha}\\
&=\sum_{\alpha\in\widetilde{\mathcal L}}\chi^{-1}(\alpha)\langle f,g[\alpha^{-1}]_k\rangle_{\alpha\Gamma_1(N)\alpha^{-1}\cap\Gamma_1(N)}\\
&=\sum_{\alpha\in\widetilde{\mathcal L}}\chi^{-1}(\alpha)\overline{\chi(\alpha^*)}\langle f,g[\alpha^*]_k^\chi\rangle_\Gamma\\
&=\overline\chi(\det\alpha)\langle f,g[T(l,m)]_k^\chi\rangle_\Gamma\\
&=\overline\chi(lm)\langle f,g[T(l,m)]_k^\chi\rangle_{\Gamma(1)}.
\end{align*}
注意 $\chi^{-1}=\overline\chi$.\footnote{译注：原书页40计算末项的下标写为 $\Gamma(1)$，而命题及首项为 $\Gamma_1(N)$；此外前文带撇的归一化在此计算中未统一写出.此处保留原式.}

\end{proof}

新的困难在于，$p\mid N$ 时 $T(p)$ 一般并不正规，因此没有一般理由使它们可对角化.故只能把 $(n,N)=1$ 的 $T(n)$ 同时对角化.

\begin{corollary*}

$\mathcal S_k(N,\chi)$ 有一组正交基，其元素为所有 $(n,N)=1$ 的 $T(n)$（乃至所有 $(lm,N)=1$ 的 $T(l,m)$）的共同特征向量.若 $f=\sum_{n>0}a(n)q^n$ 的特征值为 $(\lambda(n))_{(n,N)=1}$，则 $a(n)=a(1)\lambda(n)$，其中 $(n,N)=1$.特别地，归一化为 $a(1)=1$ 时，
\begin{align*}
a(nm)&=a(n)a(m)\quad((n,m)=1,\ (n,N)=1),\\
a(p^{r+1})&=a(p)a(p^r)-\chi(p)p^{k-1}a(p^{r-1})\quad(p\nmid N\text{ 为素数},\ r\ge1).
\end{align*}
\end{corollary*}

\begin{proof}

与 $\Gamma(1)$ 情形相同，仅有两点差别.第一式用 $a(nm)=a(m)\lambda(n)$，这由 $f[T(n)]_k^\chi$ 的Fourier展开中 $q^m$ 项的公式得出.第二式所需的 $[T(p,p)]_k^\chi$ 的特征值现在为 $\chi(p)p^{k-2}$.

\end{proof}

此时出现级为 $\Gamma(1)$ 时没有的两个问题.

\paragraph{问题.}对上述特征形式，只得到 $L$ 函数的部分分解：
\begin{align*}
L(s,f)=\prod_{p\nmid N}\frac1{1-a(p)p^{-s}+\chi(p)p^{k-1-2s}}\left(\sum_{n'}a(n')n'^{-s}\right),
\end{align*}
其中 $n'$ 遍历所有素因子都整除 $N$ 的整数.这还不令人满意，必须进一步研究.

\paragraph{问题.}这样的 $f$ 完全可能满足 $a(1)=0$！此时对所有 $(n,N)=1$ 也都有 $a(n)=0$，因此关于 $f$ 的 $L$ 函数并未得到任何信息.

\subsection{旧形式、新形式与本原形式}\label{sec:modular-1-3-6}

事实上，这两个问题彼此相关.固定特征值系 $\lambda'=(\lambda(n))_{(n,N)=1}$，记对应的特征子空间为 $\mathcal S_k(N,\chi)[\lambda']$.由于 $p\mid N$ 的 $T(p)$ 与 $T(n)$ 交换，它们保持这个空间.如果它恰好为一维，则每个 $T(p)$ 都按标量 $\lambda(p)$ 作用，故这个空间对所有 $\lambda(n)$（$n\in\mathbb N$）都是特征空间.此时 $L$ 函数确实有预期形式的Euler乘积：
\begin{align*}
L(s,f)&=\prod_{p\nmid N}\frac1{1-a(p)p^{-s}+\chi(p)p^{k-1-2s}}\prod_{p\mid N}\frac1{1-a(p)p^{-s}}\\
&=\prod_p\frac1{1-a(p)p^{-s}+\chi(p)p^{k-1-2s}}.
\end{align*}
回忆 $p\mid N$ 时 $\chi(p)=0$.如果该空间维数大于1，则事先没有理由使 $p\mid N$ 的 $T(p)$ 可对角化.不过，此时线性形式 $f\mapsto a_1(f)=a(1)$ 的核非零，即该空间包含非零形式，其所有与 $N$ 互素的指标 $n$ 都满足 $a(n)=0$.下述结果解释了这些形式的来源.

\begin{theorem*}

i) 设 $l>1$ 为 $N$ 的因子，$f\in\mathcal S_k(\Gamma_1(N/l))$.则 $i_lf:z\mapsto f(lz)$ 属于 $\mathcal S_k(\Gamma_1(N))$，且 $(n,N)=1$ 时 $a_n(i_lf)=0$.

ii) 反之，若 $f\in\mathcal S_k(\Gamma_1(N))$ 满足 $(n,N)=1$ 时 $a_n(f)=0$，则对 $N$ 的因子 $l>1$，存在 $f_l\in\mathcal S_k(\Gamma_1(N/l))$，使 $f=\sum_{l\mid N}i_lf_l$.

\end{theorem*}

\begin{proof}

i) 注意 $f(lz)=l^{1-k}f[\left[\begin{smallmatrix}l&0\\0&1\end{smallmatrix}\right]]'_k(z)$.它是群 $\left[\begin{smallmatrix}l&0\\0&1\end{smallmatrix}\right]^{-1}\Gamma_1(N/l)\left[\begin{smallmatrix}l&0\\0&1\end{smallmatrix}\right]\cap\Gamma(1)$ 的权 $k$ 自守函数，而该群包含 $\Gamma_1(N)$.它显然全纯且在尖点消失，故属于所述尖点形式空间.又其展开为 $f(lz)=\sum_{n\in\mathbb N}a_n(f)q^{ln}$，所以 $l\nmid n$ 时 $a_n(i_lf)=0$，特别地 $(n,N)=1$ 时如此.

ii) 逆命题较难，将作为一份习题课讲义的内容.

\end{proof}

这一结果说明，困难来自较低级的模形式.因此考虑 $\mathcal S_k(\Gamma_1(N))$ 的子空间 $\mathcal S_k(\Gamma_1(N))^{\mathrm{old}}$，它由所有 $f(lz)$ 生成，其中 $f\in\mathcal S_k(\Gamma_1(M))$，$M\mid N$ 且 $M\ne N$，并且 $l\mid N/M$（允许 $l=1$）.这些称为\newterm{modular-oldform}{旧形式}{oldform}.虽然目前理由还不明显，但将\newterm{modular-newform}{新形式}{newform}空间 $\mathcal S_k(\Gamma_1(N))^{\mathrm{new}}$ 定义为旧形式空间关于Petersson内积的正交补是很合适的.

\begin{lemma*}

i) 旧形式和新形式空间在 $\Gamma_0(N)$ 的权 $k$ 作用下稳定，特别地分解为
\begin{align*}
\mathcal S_k(\Gamma_1(N))^{\mathrm{old}}&=\bigoplus_{\chi\in\widehat{(\mathbb Z/N\mathbb Z)^\times}}\mathcal S_k(N,\chi)^{\mathrm{old}},\\
\mathcal S_k(\Gamma_1(N))^{\mathrm{new}}&=\bigoplus_{\chi\in\widehat{(\mathbb Z/N\mathbb Z)^\times}}\mathcal S_k(N,\chi)^{\mathrm{new}}.
\end{align*}
若 $m_\chi$ 为 $\chi$ 的导子，还有
\begin{align*}
\mathcal S_k(N,\chi)^{\mathrm{old}}&=\sum_{\substack{m_\chi\mid M\mid N,\ M\ne N\\l\mid(N/M)}}i_l(\mathcal S_k(M,\chi))
=\sum_{\substack{m_\chi\mid M\mid N,\ M\ne N\\l\mid(N/M)}}i_l(\mathcal S_k(M,\chi)^{\mathrm{new}}).
\end{align*}
ii) 设 $m_\chi\mid M\mid N$，$l$ 是 $N/M$ 的素因子或为1.则对素数 $p$，
\begin{align*}
[T(p)]_k^{\chi,N}\circ i_l=\begin{cases}i_l\circ[T(p)]_k^{\chi,M},&l\ne p,\\i_1,&l=p.\end{cases}
\end{align*}
iii) 映射 $f\mapsto f[w_N]_k$，其中 $w_N=\left[\begin{smallmatrix}0&-1\\N&0\end{smallmatrix}\right]$，诱导 $\mathcal S_k(\Gamma_1(N))$ 的一个对合，将 $\mathcal S_k(N,\chi)$ 映入 $\mathcal S_k(N,\overline\chi)$.并且
\begin{align*}
[w_N]_k\circ i_l=l^{-1}(i_1\circ[w_{N/l}]_k),
\end{align*}
所以 $[w_N]_k$ 保持旧形式空间.

\end{lemma*}

\begin{proof}

i) 设 $f\in\mathcal S_k(\Gamma_1(N/l))$，$\left[\begin{smallmatrix}a&b\\c&d\end{smallmatrix}\right]\in\Gamma_0(N)$（故 $l\mid c$）.有
\begin{align*}
(i_lf)[\left[\begin{smallmatrix}a&b\\c&d\end{smallmatrix}\right]]_k
&=l^{1-k}f[\left[\begin{smallmatrix}l&0\\0&1\end{smallmatrix}\right]]_k[\left[\begin{smallmatrix}a&b\\c&d\end{smallmatrix}\right]]_k\\
&=l^{1-k}f[\left[\begin{smallmatrix}a&bl\\c/l&d\end{smallmatrix}\right]]_k[\left[\begin{smallmatrix}l&0\\0&1\end{smallmatrix}\right]]_k
=i_l(f[\left[\begin{smallmatrix}a&bl\\c/l&d\end{smallmatrix}\right]]_k).
\end{align*}
因为 $\left[\begin{smallmatrix}a&bl\\c/l&d\end{smallmatrix}\right]\in\Gamma_0(N/l)$，括号中的形式仍属于 $\mathcal S_k(\Gamma_1(N/l))$.故旧形式空间稳定.利用 $\alpha\in\Gamma(1)$ 时 $[\alpha]_k$ 与 $[\alpha^{-1}]_k$ 的伴随关系（见前面的引理），其正交补也稳定.\footnote{译注：页43此行原书写无撇的 $[\cdot]_k$，却沿用前面带撇归一化的因子 $l^{1-k}$.此处保留.}

此外，若 $\chi$ 是模 $M$ 的Dirichlet特征标，$l\mid N/M$，上述公式还表明 $i_l(\mathcal S_k(M,\chi))\subset\mathcal S_k(N,\chi)$，从而得到所述等式的 $\supseteq$.这也给出反向包含：若 $f\in\mathcal S_k(N,\chi)^{\mathrm{old}}$ 写成 $\sum_j i_{l_j}f_j$，其中 $f_j\in\mathcal S_k(M_j,\chi_j)$，$l_j\mid N/M_j$，则也有 $f=\sum_{j,\chi_j=\chi}i_{l_j}f_j$.

ii) 对 $f\in\mathcal S_k(M,\chi)$，$q$ 展开公式给出
\begin{align*}
a_n(i_l(f[T(p)]_k^{\chi,M}))&=\begin{cases}0,&l\nmid n,\\\displaystyle\sum_{d\mid(n/l,p)}\chi(d)d^{k-1}a_{np/(ld^2)}(f),&l\mid n,\end{cases}\\
a_n((i_lf)[T(p)]_k^{\chi,N})&=\sum_{d\mid(n,p)}\chi(d)d^{k-1}a_{np/d^2}(i_lf),\\
a_{np/d^2}(i_lf)&=\begin{cases}0,&l\nmid np/d^2,\\a_{np/(ld^2)}(f),&l\mid np/d^2.\end{cases}
\end{align*}
先设 $l\ne p$.则 $l\mid np/d^2\Longleftrightarrow l\mid n$，且 $(n,p)=(n/l,p)$，于是两种计算所得 $a_n$ 相同.再设 $l=p$，于是 $p\mid N$，$\chi(p)=0$，上式化为
\begin{align*}
a_n((i_pf)[T(p)]_k^{\chi,N})=a_{np}(i_p(f))=a_p(f)=a_p(i_1f).
\end{align*}
\footnote{译注：原书页43随后两项确为 $a_p$，未写 $a_n$.此处保留.}

iii) 直接计算得
\begin{align*}
w_N\begin{bmatrix}a&b\\c&d\end{bmatrix}w_N^{-1}=\begin{bmatrix}d&-c/N\\-bN&a\end{bmatrix}.
\end{align*}
因此共轭给出 $\Gamma_1(N)$ 和 $\Gamma_0(N)$ 的对合.映射 $f\mapsto f[w_N]'_k$ 因而诱导权 $k$、级 $\Gamma_1(N)$ 的自守函数的对合，显然保持模形式和尖点形式.共轭还诱导特征标的对合 $\chi\mapsto{}^{w_N}\chi$，且
\begin{align*}
{}^{w_N}\chi\left(\begin{bmatrix}a&b\\c&d\end{bmatrix}\right)
=\chi\left(w_N\begin{bmatrix}a&b\\c&d\end{bmatrix}w_N^{-1}\right)=\chi(a)=\overline\chi(d)=\overline\chi\left(\begin{bmatrix}a&b\\c&d\end{bmatrix}\right).
\end{align*}
故 ${}^{w_N}\chi=\overline\chi$，从而 $f\mapsto f[w_N]_k$ 把相应特征标空间映到共轭特征标空间.最后由显式公式 $f[w_N]_k(z)=N^{-1}z^{-k}f(-1/(Nz))$ 计算复合即可.\footnote{译注：原书页43这里先写带撇的作用，末句显式公式却写无撇并取 $N^{-1}$；页42同时称该映射为对合且写因子 $l^{-1}$.不同归一化之间有疑问，均照原式保留.}

\end{proof}

\begin{remark*}

由 i)，若 $\chi$ 本原，即 $m_\chi=N$，则 $\mathcal S_k(N,\chi)=\mathcal S_k(N,\chi)^{\mathrm{new}}$.

\end{remark*}

\begin{corollary*}

$\mathcal S_k(N,\chi)^{\mathrm{old}}$ 与 $\mathcal S_k(N,\chi)^{\mathrm{new}}$ 在 $\mathcal H_0(N)$ 的作用下稳定.

\end{corollary*}

\begin{proof}

由引理 ii)，旧形式空间在所有 $T(p)$ 下稳定；又显然在 $p\nmid N$ 的 $T(p,p)$ 下稳定，故在整个Hecke环下稳定.因此新形式空间在 $T(p)$ 的伴随算子下稳定.$p\nmid N$ 时，伴随为 $\overline\chi(p)T(p)$，所以 $T(p)$ 保持新形式空间.$p\mid N$ 时须另作论证！

此时 $\mathcal L_0(N)_{1,p}=\{\left[\begin{smallmatrix}1&j\\0&p\end{smallmatrix}\right]:0\le j<p\}$，故 $f[T(p)]_k^\chi=\sum_{j=0}^{p-1}f[\left[\begin{smallmatrix}1&j\\0&p\end{smallmatrix}\right]]_k$.用于证明 $p\nmid N$ 时正规性的计算，对每个 $p$ 都给出伴随算子
\begin{align*}
T(p)^*:f\longmapsto\sum_{j=0}^{p-1}f[\left[\begin{smallmatrix}p&j\\0&1\end{smallmatrix}\right]]_k.
\end{align*}
此算子正是 $w_N\circ[T(p)]_k^{\overline\chi}\circ w_N^{-1}$.由引理 ii)、iii)，它保持旧形式空间；取伴随便知 $T(p)$ 保持新形式空间.\footnote{译注：页44顶部的伴随矩阵上右元确写 $j$，这里未据矩阵求逆改为其他位置或符号.}

\end{proof}

现在可运用基于部分特征值系重数1的论证了.

\begin{corollary*}

新形式空间 $\mathcal S_k(N,\chi)^{\mathrm{new}}$ 有一组正交基，由所有 $T(n)$（$n\in\mathbb N$）的共同特征形式组成.每个特征值系的重数为1.

\end{corollary*}

\begin{proof}

由前述论证，只须证明部分特征值系 $\lambda'=(\lambda(n))_{(n,N)=1}$ 的重数为1，即相应特征空间为一维.否则该空间包含非零形式 $f$ 满足 $a_1(f)=0$，因而 $(n,N)=1$ 时 $a_n(f)=0$.但前述定理 ii) 说明这样的 $f$ 是旧形式，矛盾.

\end{proof}

新、特征且归一化的形式 $f\in\mathcal S_k(\Gamma_1(N))$ 称为导子为 $N$ 的\newterm{modular-primitive-form}{本原形式}{primitive form}.本原形式组成新形式空间的一组基.引理 i) 表明，级为 $\Gamma_1(N)$ 的所有形式都由级为 $\Gamma_1(M)$（$M\mid N$）的本原形式构造而来.实际上，可以证明下述定理（此处不证）.

\begin{theorem*}

集合 $\{f(lz):f\in\mathcal S_k(\Gamma_1(M))\text{ 为本原形式，且 }Ml\mid N\}$ 是 $\mathcal S_k(\Gamma_1(N))$ 的一组基.

\end{theorem*}

